% answers_short.tex —— 1992 年·全国卷·理科 · 简短答案速查
% 由 tools/make_short_answers.py 自 answers.tex 的 \gksolhead/\gkanswer 提取，
% 并与 1992/1992_数学真题_全国理_答案要点.md 逐题比对后生成，勿手改（改源 MD 后重跑）。

\begin{gknotice}
\textbf{使用说明}：本节只给最终答案便于快速核对，完整推导见解析版。答案为编写者独立推导，\textbf{非官方标准答案}。
\end{gknotice}

\gkpart{选择题}
\begin{center}\begin{tabular}{ccccccccccccccccccc}
\toprule
题号 & 1 & 2 & 3 & 4 & 5 & 6 & 7 & 8 & 9 & 10 & 11 & 12 & 13 & 14 & 15 & 16 & 17 & 18 \\
\midrule
答案 & A & D & D & B & D & B & B & C & D & D & B & B & A & D & D & C & A & C \\
\bottomrule
\end{tabular}\end{center}

\gkpart{填空题}
\begin{itemize}[leftmargin=2.6em]
  \item[\textbf{19.}] \(x=-1\)
  \item[\textbf{20.}] \(\dfrac14\)
  \item[\textbf{21.}] \(\dfrac{15}{128}\)
  \item[\textbf{22.}] \(\dfrac{(x-2)^2}{4}-\dfrac{y^2}{12}=1\)
  \item[\textbf{23.}] \(\dfrac{13}{16}\)
\end{itemize}

\gkpart{解答题}
\begin{itemize}[leftmargin=2.6em]
  \item[\textbf{24.}] \(z=-1\) 或 \(z=-1+3\mathrm{i}\)。
  \item[\textbf{25.}] \(\sin 2\alpha=-\dfrac{56}{65}\)。
  \item[\textbf{26.}] 成立（\(\overrightarrow{EF}^{\,2}=d^{2}+m^{2}+n^{2}\pm2mn\cos\theta\)，故 \(EF=\sqrt{d^{2}+m^{2}+n^{2}\pm2mn\cos\theta}\)）。
  \item[\textbf{27.}] (1) \(-\dfrac{24}{7}<d<-3\)；(2) \(S_6\) 最大
  \item[\textbf{28.}] 成立（由中点弦斜率关系得 \(x_0=\dfrac{a^{2}-b^{2}}{a^{2}}x_m\)，结合 \(|x_m|<a\) 即得结论）。
\end{itemize}
